\documentclass[11pt,a4paper]{article}
\usepackage[utf8]{inputenc}
\usepackage[T1]{fontenc}
\usepackage[ngerman]{babel}
\usepackage{lmodern}
\usepackage[margin=25mm]{geometry}
\usepackage{url}
\usepackage{enumitem}
\urlstyle{tt}
\setlist[description]{style=nextline,leftmargin=1em,itemsep=2pt,parsep=0pt}
\newcommand{\datei}[1]{\path{#1}}
\setlength{\parindent}{0pt}
\setlength{\parskip}{6pt}
\setlength{\emergencystretch}{2em}

\begin{document}

\section{runStaticStudy}

\subsection{Aufgabe des Hauptskripts}
Das Hauptskript dient der Koordination von Datenerzeugung, Modellkonfiguration
und -training, Auswertung sowie Laufzeitvergleich und Berichtserstellung.
Über das Argument \texttt{phase} wird festgelegt, welcher Verarbeitungsschritt
ausgeführt werden soll. Die Argumente \texttt{configFile} und
\texttt{localFile} bestimmen die fachlichen Einstellungen sowie die
verwendeten Datenpfade und den Ausgabeordner. Die eigentlichen Berechnungen
übernehmen die jeweils aufgerufenen Funktionen.

\subsection{Aufruf und Ablauf}
\begin{verbatim}
output = runStaticStudy(phase, configFile, localFile)
\end{verbatim}
% Hier den Ablauf und die Rueckgabe in eigenen Worten beschreiben.
\vspace{2\baselineskip}

\subsection{Argument: phase}
% Hier die einzelnen Phasen nach und nach erklaeren.
% help, init, generate, tune, train, freeze, evaluate, benchmark, report
% Bei Bedarf je Phase einen Unterabschnitt mit \subsubsection{...} anlegen.
\vspace{2\baselineskip}

\subsection{Argument: configFile}
% Hier meine Erklaerung und eigene Beispiele ergaenzen.
\vspace{2\baselineskip}

\subsection{Argument: localFile}
% Hier meine Erklaerung und eigene Beispiele ergaenzen.
\vspace{2\baselineskip}

\section{FEM-Datenerzeugung}

\subsection{Ein einzelner FEM-Fall}
Datei: \texttt{fem\_core/src/solve\_case.m}

Die Funktion \texttt{solve\_case} berechnet für eine vorgegebene
Parameterkombination einen einzelnen statischen FEM-Fall. Dazu erstellt sie
ein dreidimensionales, linear-elastisches Modell, löst es und bereitet das
berechnete Verschiebungsfeld für die weitere Verwendung auf.

\subsection{Eingaben und Berechnung}
% Hier Parameter, Geometrie, Material, Einspannung, Last und Loesung erklaeren.
\vspace{2\baselineskip}

\subsection{Ergebnis}
% Hier das Ergebnis und seine Weiterverwendung beschreiben.
\vspace{2\baselineskip}

\newpage
\appendix
\raggedright
\section{Bestandteile des Projekts}

Dieser Anhang gibt eine Übersicht über den aktiven statischen Studienablauf.
Alle Pfade beziehen sich auf \textbf{Code v2}. Nicht jede Datei muss für den
Einstieg im Detail verstanden werden. Viele Ergebnisse entstehen mehrfach;
dafür sind unten die Dateinamen und Ordner als Muster angegeben.

\subsection{Fünf Stationen zum Einstieg}
\begin{enumerate}[nosep]
\item \datei{runStaticStudy.m}: Welche Phase wird ausgeführt?
\item \datei{study.json} und \datei{local.example.json}: Welche Einstellungen und Daten werden verwendet?
\item \datei{generate.m} und \datei{solve_case.m}: Wie entstehen FEM-Daten?
\item \datei{fitModel.m} und eine Modellmethode: Wie wird ein Modell trainiert?
\item \datei{evaluation.m} und \datei{benchmark.m}: Wie werden Genauigkeit und Laufzeit untersucht?
\end{enumerate}

\subsection{Voraussetzungen, Einstieg und Konfiguration}
Für den vollständigen Ablauf werden MATLAB mit Deep Learning Toolbox,
Statistics and Machine Learning Toolbox und Partial Differential Equation
Toolbox sowie Java, Git und Python mit Matplotlib benötigt. Der Studienlauf
prüft während der Ausführung eine Speicherplatzreserve von 512 MiB.
\begin{description}
\item[\datei{runStaticStudy.m}] Startet jeweils eine Phase der Studie.
\item[\datei{setupProject.m}] Fügt Studien- und FEM-Ordner zum MATLAB-Suchpfad hinzu.
\item[\datei{workflows/static_comparison/configs/study.json}] Fachliche Einstellungen, Methoden, Kandidaten und Trainingsgrößen.
\item[\datei{workflows/static_comparison/configs/local.example.json}] Lokale Datenpfade, Python und Ausgabeordner; für einen neuen Lauf eine Kopie mit neuem Ausgabeordner verwenden.
\item[\datei{workflows/static_comparison/configs/training_blocks.csv}] Zuordnung der Trainingsfälle zu den beiden Trainingsblöcken.
\item[\datei{workflows/static_comparison/studyPlots.py}] Erstellt die Diagramme der Berichtsphase.
\end{description}

\subsection{Vorhandene Eingabedaten}
\begin{description}
\item[\datei{data_generation/data/static/train.mat}] Vorhandener Trainingspool mit FEM-Verschiebungsfeldern.
\item[\datei{data_generation/data/static/design.csv}] Fall-IDs und Parameter; wird unter anderem für Überschneidungsprüfungen verwendet.
\item[\datei{training_blocks.csv}] Metadaten zu den Trainingsblöcken im Konfigurationsordner.
\end{description}
Bei \texttt{reference\_mode="reuse"} werden außerdem Design und Referenzdaten
aus dem eingestellten vorherigen Lauf verwendet: \datei{candidate_design.mat},
\datei{candidate_design.csv}, \datei{candidate_design.sha256},
\datei{reference_costs.csv} sowie unter \datei{references/val/},
\datei{references/interpolation/} und \datei{references/extrapolation/} die
Dateien \datei{val.mat} beziehungsweise \datei{test.mat}, jeweils mit
\datei{design.csv}, \datei{data.sha256}, \datei{ledger.csv} und
\datei{generation_sessions.csv}. Für einen anwendbaren Vergleich mit dem
früheren Lauf dienen dort auch \datei{run_results.csv} und \datei{selections.mat}.
Bei \texttt{reference\_mode="generate"} werden neue Validierungs- und
Testreferenzen erzeugt; der Trainingspool wird separat bereitgestellt.

\subsection{FEM und gemeinsame Schnittstellen}
\begin{description}
\item[\datei{fem_core/staticFemConfig.m}] Stellt die FEM-Einstellungen zusammen.
\item[\datei{fem_core/src/solve_case.m}] Berechnet einen einzelnen statischen FEM-Fall.
\item[\datei{fem_core/src/reference_grid.m}] Erstellt Auswertegitter und Gewichte.
\item[\datei{loadFemDataset.m}] Lädt FEM-Daten im zugrunde liegenden Format.
\end{description}
Im Paket \datei{+surrogate/} liegen die gemeinsamen Funktionen:
\begin{description}
\item[\datei{contract.m}] Definiert Parameter, Einheiten und Auswertegitter.
\item[\datei{loadDataset.m}, \datei{validateDataset.m}] Laden und prüfen Datensätze.
\item[\datei{validateTrainingSplits.m}] Prüft Trainings- und Validierungsdaten.
\item[\datei{fitModel.m}, \datei{predictModel.m}] Gemeinsame Aufrufe für Training und Vorhersage.
\item[\datei{evaluatePrediction.m}] Vergleicht Vorhersagen mit Referenzen.
\item[\datei{saveModel.m}, \datei{loadModel.m}] Speichern und laden Modelle.
\end{description}
Unter \datei{+surrogate/+study/} prüfen
\datei{assertDisjointParameters.m} die Trennung von Parametern,
\datei{nestedTrainingSubsets.m} erstellt Trainingsauswahlen und
\datei{summarizeErrors.m} fasst Fehler zusammen.

\subsection{Die vier Modellverfahren}
Unter \datei{+surrogate/+models/} besitzen die vier Methoden jeweils eine
\datei{fit.m} für das Training und eine \datei{predict.m} für die Vorhersage:
\begin{description}
\item[\datei{+mlp/}] Direktes neuronales Netz.
\item[\datei{+rbf/}] Netz mit RBF-Merkmalen.
\item[\datei{+pod_mlp/}] POD mit einem neuronalen Netz für die Koeffizienten.
\item[\datei{+pod_gpr/}] POD mit Gaußprozessregression für die Koeffizienten.
\end{description}
Der Ordner \datei{+shared/} enthält \datei{fitNeural.m} und
\datei{predictNeural.m} für MLP und RBF, \datei{rbfFeatures.m} für die
RBF-Merkmale, \datei{fitPod.m} für die POD-Basis sowie
\datei{podFingerprint.m} und \datei{podDiagnostics.m} für Prüfung und Diagnostik.

\subsection{Funktionen zur Studiensteuerung}
Alle folgenden Dateien liegen unter \datei{workflows/static_comparison/+staticstudy/}.
\begin{description}
\item[\datei{configure.m}, \datei{initialize.m}] Lesen Einstellungen und bereiten den Lauf vor.
\item[\datei{createDesign.m}] Erstellt das Design für neue Referenzfälle.
\item[\datei{sources.m}, \datei{hash.m}, \datei{verify.m}] Erfassen und prüfen Quellen und Eingaben.
\item[\datei{requireSpace.m}, \datei{requireBudget.m}] Prüfen Speicherplatz und ein gegebenenfalls gesetztes Zeitbudget.
\item[\datei{generate.m}, \datei{recordGeneration.m}] Erzeugen FEM-Referenzen und protokollieren deren Laufzeit.
\item[\datei{validateReferences.m}, \datei{loadReference.m}, \datei{reference.m}] Prüfen und laden Referenzen; regeln den Zugriff auf Testdaten.
\item[\datei{candidates.m}, \datei{tuning.m}, \datei{recordSearch.m}] Erstellen und untersuchen Modellkandidaten und protokollieren die Suche.
\item[\datei{trainingBlocks.m}, \datei{subset.m}, \datei{distribution.m}] Stellen Trainingsauswahlen und Verteilungen zusammen.
\item[\datei{fitCheckpoint.m}, \datei{training.m}] Trainieren und speichern einzelne Fits beziehungsweise organisieren die Studienläufe.
\item[\datei{freeze.m}, \datei{checkFrozen.m}] Halten den Stand fest und prüfen ihn.
\item[\datei{evaluation.m}, \datei{observeModel.m}, \datei{timePrediction.m}] Werten Modelle aus und erfassen Vorhersagezeiten.
\item[\datei{benchmark.m}] Vergleicht FEM- und Modelllaufzeiten.
\item[\datei{report.m}, \datei{comparePrevious.m}] Erstellen Zusammenfassungen und gegebenenfalls einen Vergleich mit dem vorherigen Lauf.
\end{description}

\subsection{Erzeugte Dateien: Initialisierung und Referenzen}
Die folgenden Pfade liegen im Ausgabeordner des jeweiligen Studienlaufs.
\begin{description}
\item[\datei{resolved_config.json}] Tatsächlich verwendete Konfiguration.
\item[\datei{source_manifest.csv}, \datei{environment.mat}] Quell- und Eingabestand sowie Angaben zur Umgebung.
\item[\datei{reference_locations.json}] Orte der Validierungs- und Testreferenzen.
\item[\datei{candidate_design.mat}, \datei{candidate_design.csv}, \datei{candidate_design.sha256}] Design und Prüfsumme; neu erstellt oder übernommen.
\end{description}
Bei neuer Referenzerzeugung entstehen unter \datei{references/val/},
\datei{references/interpolation/} und \datei{references/extrapolation/}:
\begin{description}
\item[\datei{cases/case_<ID>.mat} und \datei{case_<ID>.mat.sha256}] Einzelner FEM-Fall und Prüfsumme.
\item[\datei{val.mat} oder \datei{test.mat}] Zusammengesetzter Datensatz.
\item[\datei{design.csv}, \datei{ledger.csv}] Ausgewählte Eingaben und Protokoll der Berechnungsversuche.
\item[\datei{generation_sessions.csv}, \datei{data.sha256}] Zeitprotokoll und Prüfsumme des Datensatzes.
\end{description}
Bei Wiederverwendung bleiben die Referenzen im eingestellten vorhandenen Ordner.

\subsection{Erzeugte Dateien: Tuning, Training und Freeze}
Unter \datei{tuning/<Methode>/candidate_<Nummer>_seed<Seed>/} und in den
Modellordnern des Trainings stehen typischerweise \datei{model.mat},
\datei{result.mat} und \datei{result.sha256}: Modell, Fit-Ergebnis und
Prüfsumme. Fehlgeschlagene Fits besitzen nicht notwendigerweise ein Modell.
\begin{description}
\item[\datei{tuning/<Methode>/candidate_summary.csv}] Fehlerübersicht der Kandidaten.
\item[\datei{tuning/<Methode>/search_summary.mat}, \datei{search_sessions.csv}] Suchergebnis und Zeitprotokoll je Methode.
\item[\datei{search_sessions.csv}] Übergreifendes Zeitprotokoll im Ausgabeordner.
\item[\datei{selections.mat}, \datei{selections.sha256}] Gewählte Einstellungen und Prüfsumme.
\item[\datei{runs/}, \datei{controlled/}, \datei{distribution/}, \datei{timing_repeats/}] Hauptläufe, Rangvergleiche, Verteilungsvergleiche und Zeitwiederholungen.
\item[\datei{bases/}] Enthält \datei{pod.mat} und \datei{pod.sha256} für gespeicherte POD-Basen.
\item[\datei{distribution_ids.mat}, \datei{training.complete}] Fallauswahlen und Abschlusskennzeichen des Trainings.
\item[\datei{freeze_manifest.csv}, \datei{freeze_manifest.sha256}] Festgeschriebene Dateien und Prüfsumme des Manifests.
\end{description}
Bei Unterbrechungen oder ungültigen Rängen können zusätzliche Dateien wie
\datei{pod_incomplete.mat}, \datei{invalid_rank.txt} oder
\datei{case_<ID>.mat.incomplete_...} entstehen.

\subsection{Erzeugte Dateien: Auswertung, Benchmark und Bericht}
Unter \datei{evaluation/} entstehen für jeden ausgewerteten Modelllauf
\datei{interpolation.mat} und \datei{extrapolation.mat}, jeweils mit einer
\datei{.sha256}-Datei. \datei{evaluation.complete} kennzeichnet den Abschluss.
\begin{description}
\item[\datei{benchmark/benchmark_cases.csv}, \datei{benchmark_meta.mat}] Fälle und Einstellungen des Laufzeitvergleichs.
\item[\datei{benchmark/fem_partial.mat}, \datei{surrogate_partial.mat}] Zwischenstände der Zeitmessungen.
\item[\datei{benchmark/benchmark_results.mat} und \datei{benchmark_results.mat.sha256}] Gesamte Benchmarkdaten und Prüfsumme.
\item[\datei{benchmark/fem_raw.csv}, \datei{surrogate_raw.csv}, \datei{batch_raw.csv}] Einzelne Zeitmessungen.
\item[\datei{benchmark/summary.csv}, \datei{break_even.csv}] Zusammenfassung und Break-even-Szenarien.
\item[\datei{benchmark/BENCHMARK.md}, \datei{benchmark.complete}] Beschreibung und Abschlusskennzeichen.
\item[\datei{run_results.csv}, \datei{case_results.csv}, \datei{pipeline_summary.csv}] Ergebnisse je Lauf, Fall und zusammengefasster Studiengruppe.
\item[\datei{search_costs.csv}, \datei{search_attempts.csv}, \datei{gpr_fit_timing_repeats.csv}] Suchkosten, Kandidatenversuche und GPR-Zeitwiederholungen.
\item[\datei{reference_costs_historical.csv}, \datei{REPORT.md}, \datei{report.complete}] Historische Referenzkosten, Bericht und Abschlusskennzeichen.
\item[\datei{comparison_previous.csv}, \datei{VERGLEICH.md}] Gegebenenfalls Vergleich mit dem vorherigen Lauf.
\item[\datei{plots/*.png}] Diagramme aus der Berichtsphase.
\end{description}
Bei einer unterbrochenen Auswertung können zusätzliche Dateien mit
\datei{.incomplete} entstehen. Weitere Abschlussprüfungen eines konkreten
Laufs können zusätzliche Dokumente erzeugen; sie gehören nicht automatisch
zu jeder Studie.

\subsection{Beispiele und Tests}
Die Beispiele \datei{examples/trainStaticModelExample.m} und
\datei{examples/showStaticResultsExample.m} zeigen ein einzelnes Training
beziehungsweise vorhandene Ergebnisse. Sie sind für die Hauptstudie nicht
erforderlich.

Im Ordner \datei{tests/} prüfen \datei{runStaticValidation.m} und
\datei{runStaticWorkflowSmoke.m} das Projekt beziehungsweise einen kleinen
vollständigen Studienablauf. \datei{staticWorkflowTest.m},
\datei{directNeuralTest.m}, \datei{podMlpTest.m}, \datei{podGprTest.m},
\datei{surrogateContractTest.m}, \datei{nestedTrainingSubsetsTest.m} und
\datei{weightedDirectLossTest.m} enthalten die zugehörigen Unit-Tests;
\datei{makeContractData.m} stellt kleine Testdaten bereit. Unter
\datei{tests/fixtures/} liegen \datei{study_baseline.mat} und
\datei{regression.json} für Regressionsprüfungen. Der Smoke-Test verwendet
\datei{configs/smoke.json} und eine eigene lokale Konfiguration; er erzeugt
einen separaten Testlauf.

\end{document}
